\section{Introduction}
\label{sec:intro}

Three problems from three fields have the same shape.

\paragraph{A cast.} In the proof assistant Lean~4 \citep{demoura2021}, the two expressions
below differ only in where a natural number is turned into an integer. The first evaluates
to $0$, the second to $-1$:
\begin{lstlisting}
example : ((2 - 3 : Nat) : Int) = 0 := rfl
example : ((2 : Nat) - (3 : Nat) : Int) = -1 := rfl
\end{lstlisting}
The symbol \code{-} is the same in both lines. What it means depends on which side of the
cast it is evaluated on, and the author of the expression does not write the cast.

\paragraph{A verified kernel.} The functional-correctness proof of the seL4 microkernel
establishes that its C implementation refines an abstract specification
\citep{klein2009}. The proof has no notion of time and so makes no claim about timing
channels. Isolation in time had to be added later, as a mechanism of its own
\citep{ge2019}. The kernel was verified, and a property a user might take ``verified'' to
include was not among the properties the proof carried from the specification to the
code.

\paragraph{A valid chain of reasoning.} A language model writes out a chain of steps
before its answer, or a proof that a proof checker accepts. A check on that text can
establish that every step follows from the ones before it. It does not establish that the
text describes how the model arrived at its answer. In controlled studies, models whose
answers were changed by a hint often did not mention the hint in the reasoning they wrote
\citep{turpin2023,chen2025}.

\medskip
In each case a claim is established on one side of a boundary and reused on the other: an
identity on $\N$ reused on $\Z$, a property of a specification reused of an
implementation, the validity of a text reused as an account of a computation. Whether the
reuse is sound depends on the map or relation that crosses the boundary, on what kind of
crossing it is and on what it preserves. In each case the thing that is reused, a symbol,
the word ``verified'', the word ``reasoning'', does not say.

\paragraph{Thesis.} A crossing along which claims are carried should be declared and
typed: by its kind, which fixes what a failure of the crossing consists in, and by what it
preserves, under which conditions.

\paragraph{What the fields already do.} None of the three fields lacks this insight. Each
applies it locally. Lean's mathematical library Mathlib records what a cast preserves in
bundled morphisms and in lemmas registered for its cast tactics
\citep{mathlib2020,lewis2020}. Verification practice proves the preservation of each
further property separately, and the theory of secure compilation classifies compilers by
the class of properties they preserve \citep{clarkson2010,abate2019}. Interpretability
research separates the faithfulness of an explanation from its plausibility and tests the
former by intervening on the model \citep{jacovi2020,geiger2025}. What these local
solutions have in common is not usually stated, because the three literatures do not
share a vocabulary. This paper proposes one.

\paragraph{Kind of paper.} This is a position paper with a machine-checked artefact. It
reports no new experiments and proves no deep theorem. Its claims are of three sorts:
readings of existing practice, which can be checked against the literature cited;
definitions and elementary results, some of which are checked in Lean; and two
predictions, which are marked as untested.

\paragraph{Contributions.}
\begin{enumerate}[label=(\arabic*)]
  \item \emph{A common reading of three practices.} Cast lemmas in Mathlib, per-property
        preservation proofs in system verification, and faithfulness tests for written
        model reasoning are shown to address one situation: a claim established on one
        side of a crossing is reused on the other
        (Sections~\ref{sec:casts}--\ref{sec:reasoning}).
  \item \emph{A small vocabulary for declaring a crossing.} A crossing is declared by its
        kind and by what it preserves (Section~\ref{sec:rule}). The kind cannot be
        inferred from the map alone, because two maps with the same formal properties can
        differ in the role in which they are used (Section~\ref{sec:three-kinds}). The
        vocabulary is taken from Layered Event Modelling \citep{kreinzlem} and restated
        here in the minimal form the argument needs.
  \item \emph{What a check on a trace certifies.} The validity of a written trace and its
        faithfulness to the computation that produced it are crossings of different
        kinds. A judgement made on the trace alone can agree with faithfulness exactly
        when faithfulness does not vary among the computations that write the same
        trace. Where one trace is written by two computations and describes only one of
        them, no such judgement exists (Section~\ref{sec:kernel}).
  \item \emph{A partial mechanisation.} The graph model behind the vocabulary, the
        typed crossing with two worked examples from the first case, and the results on
        traces are formalised in Lean~4 and published as an open-source package
        \citep{kreinzlemlean}. The refinement case is not formalised.
        Section~\ref{sec:mech} states what the mechanisation does and does not
        establish.
  \item \emph{Two testable predictions} about where automated provers fail and about
        what a declaration of kind adds
        (Section~\ref{sec:limits}).
\end{enumerate}

\paragraph{What is not claimed.} The paper does not claim that any of the three fields
has overlooked the problem; that the vocabulary of Layered Event Modelling transfers to
other domains in general, since three cases chosen by the author are evidence and not
proof; that written model reasoning is usually unfaithful, or of no value as evidence; or
that checking the results in Lean shows the definitions to fit real systems.

\paragraph{Outline.} Section~\ref{sec:rule} states the rule and the vocabulary.
Sections~\ref{sec:casts}, \ref{sec:refinement} and~\ref{sec:reasoning} are the three
cases; each is written to be read without the others and without background in the other
fields. Section~\ref{sec:mech} describes the mechanisation, Section~\ref{sec:related} the
related work, and Section~\ref{sec:limits} the limitations and the predictions.
